\documentclass[12pt, letterpaper]{article}
\date{\today}
\usepackage[margin=1in]{geometry}
\usepackage{amsmath}
\usepackage{hyperref}
\usepackage{cancel}
\usepackage{amssymb}
\usepackage{fancyhdr}
\usepackage{pgfplots}
\usepackage{booktabs}
\usepackage{pifont}
\usepackage{amsthm,latexsym,amsfonts,graphicx,epsfig,comment}
\pgfplotsset{compat=1.16}
\usepackage{xcolor}
\usepackage{tikz}
\usetikzlibrary{shapes.geometric}
\usetikzlibrary{arrows.meta,arrows}
\newcommand{\Z}{\mathbb{Z}}
\newcommand{\N}{\mathbb{N}}
\newcommand{\R}{\mathbb{R}}
\newcommand{\Q}{\mathbb{Q}}
\newcommand{\C}{\mathbb{C}}
\newcommand{\F}{\mathbb{F}}

\newcommand{\Po}{\mathcal{P}}
\newcommand{\Pro}{\mathbb{P}}
\author{Alex Valentino}
\title{501 homework}
\pagestyle{fancy}
\renewcommand{\headrulewidth}{0pt}
\renewcommand{\footrulewidth}{0pt}
\fancyhf{}
\rhead{
	Homework 4\\
	501
}
\lhead{
	Alex Valentino\\
}
\begin{document}
\begin{enumerate}
	\item[1] Let $X$ be a set, $\mu^*$ be an outer measure on $X$,
	$\{A_j\}_{j \in \N}$ be a sequence pf $\mu^*$ measurable disjoint sets, and 
	$E \subset X$.  We want to show that
	$\mu^*(E \cap ( \bigcup_{i=1}^\infty A_j)) = \sum_{i=1}^\infty \mu^*(E \cap A_j) $.  Note that $$\mu^*(E \cap ( \bigcup_{i=1}^\infty A_j)) = 
	\mu^*(\bigcup_{i=1}^\infty (E \cap A_j)) \leq \sum_{i=1}^\infty \mu^*(E \cap A_j)
	$$ by the subadditivity of $\mu^*$.  Therefore we must show the inequality in 
	the opposite direction. If we consider that for arbitrary $j$
	\begin{align*}
	\mu^*(E \cap (\cup_{i=1}^\infty A_i)) &= \mu^*(E \cap (\cup_{i=1}^\infty A_i) \cap A_j) + \mu^*(E \cap (\cup_{i=1}^\infty A_i) \cap A_j^c)\\
	&= \mu^*(E \cap A_j) + \mu^*(E \cap (\cup_{i=1, i \neq j}^\infty A_i))
	\end{align*}

	since $(\cup_{i=1}^\infty A_i) \cap A_j = \cup_{i=1}^\infty (A_i \cap A_j) =
	A_j$ as if $i \neq j$ then $A_i \cap A_j = \emptyset$ by definition
	and $(\cup_{i=1}^\infty A_i) \cap A_j^c = \cup_{i=1}^\infty (A_i \cap A_j^c) = \cup_{i=1, i \neq j}^\infty A_i)
	$ since for all $i \neq j, A_j^c \cap A_i = A_i$.  Therefore we have by induction 
	that 
	$\mu^*(E \cap ( \bigcup_{i=1}^\infty A_j)) = \sum_{i=1}^n \mu^*(E \cap A_i) + 
	\mu^*(E \cap ( \bigcup_{i=n+1}^\infty A_j))$.  Therefore $mu^*(E \cap ( \bigcup_{i=1}^\infty A_j)) \geq \sum_{i=1}^n \mu^*(E \cap A_i)$.  Since this holds for all $n$, then 
	it holds for the limit.  Therefore we have equality.  
	\item[2] Let $\mu^*$ be an outer measure on $X$ with $\mu^*(X) < \infty$, and let 
	$\nu^*(E) = \sqrt{\mu^*(E)}$ be defined for all $E \subseteq X$.  To show that 
	$\nu^*$ is an outer measure, we must show the following:
	\begin{itemize}
		\item $\nu^*(\emptyset) = \sqrt{\mu^*(\emptyset)} = \sqrt{0} = 0$
		\item If $B_1 \subseteq B_2 \subseteq X$ then we know that $\mu^*(B_1) \leq \mu^*(B_2)$.  Therefore $\sqrt{\mu^*(B_1)} \leq \sqrt{\mu^*(B_2)}$, which implies that
		$\nu^*(B_1) \leq \nu^*(B_2)$
		\item Let $\{B_j\}_{j\in \N} \subseteq 2^X $.  Then consider that 
		\begin{align*}
			\nu^*(\cup_{j=1}^\infty B_j) &= \sqrt{\mu^*(\cup_{j=1}^\infty B_j)}\\
			&\leq \sqrt{\sum_{j=1}^\infty \mu^*( B_j)}\\
			&= \sqrt{\sum_{j=1}^\infty \nu^*( B_j)^2}\\
			&\leq \sum_{j=1}^\infty \nu^*( B_j)
		\end{align*}
		Note that since all the values are positive and infinity is an exceptable value
		that the inequalities hold.
\end{itemize}
	Now to show that $A$ is $\nu^*$ measurable iff $\nu^*(A^c) = 0$ or $\nu^*(A) = 0$.
	\begin{itemize}
		\item $(\Rightarrow)$.  Suppose $A$ is $\nu^*$ measurable.  Then 
		$\nu^*(X) = \nu^*(A) + \nu^*(A^c)$.  Therefore, 
		\begin{align*}
			\sqrt{\mu^*(X)} &= \sqrt{\mu^*(A)} + \sqrt{\mu^*(A^c)}\\
			\mu^*(X) &= \mu^*(A) + \mu^*(A^c) + 2 \sqrt{\mu^*(A)\mu^*(A^c)}\\
			&\geq \mu^*(X) + 2 \sqrt{\mu^*(A)\mu^*(A^c)}\\
			\mu^*(X) &\geq \mu^*(X) + 2 \sqrt{\mu^*(A)\mu^*(A^c)}\\
			0 \geq 2 \sqrt{\mu^*(A)\mu^*(A^c)}\\
		\end{align*}
		Note that this final inequality implies that either $\mu^*(A) = 0$ or 
		$\mu^*(A^c) = 0$.  Since these inequalities are maintained under the square root 
		then the first direction has been shown.
		\item $(\Leftarrow)$  Suppose WLOG that $\nu^*(A^c) = 0$.  
		Then for arbitrary $E \subseteq X$, 
		$0 \leq \nu^*(E \cap A^c) \leq \nu^*(A^c) = 0$ which implies that 
		$\nu^*(E \cap A^c) = 0$. Therefore 
		$\nu^*(E) \leq \nu^*(A \cap E) + \nu^*(A^c \cap E) = \nu^*(A \cap E) \leq \mu^*(E)$ which implies 
		that $\nu^*(E) = \nu^*(A \cap E) + \nu^*(A^c \cap E)$.  
\end{itemize}		 
	\item[3] Let $\mu^*$ be the Lesbegue outer measure, and $A \subset [0,1]$.  We must show that 
	$1 = \mu^*(A) + \mu^*(A^c)$ iff $A$ is Lesbesgue measurable.  
	\begin{itemize}
		\item $(\Rightarrow)$ Suppose $A$ is $\mu^*$-measurable.  Then 
		$1 = \mu^*([0,1]) = \mu^*(A) + \mu^*(A^c)$.
		\item $(\Leftarrow)$ Suppose not.  Then $1 = \mu^*(A) + \mu^*(A^c)$ and 
		$A$ is not $\mu^*$-measurable.  Therefore there exists $E \subset [0,1]$ such that 
		$\mu^*(E) \neq \mu^*(E \cap A) + \mu^*(E \cap A^c)$.  We know by subadditivity 
		that $\mu^*(E) \leq \mu^*(E \cap A) + \mu^*(E \cap A^c)$, there our inequality 
		is strict with $\mu^*(E) < \mu^*(E \cap A) + \mu^*(E \cap A^c)$.  If we assume our
		set $E$ is $\mu^*$-measurable then we have that 
		\begin{align*}
		1 &= \mu^*(A) + \mu^*(A^c)\\
		&= \mu^*(A\cap E) + \mu^*(A \cap E^c) + \mu^*(A^c\cap E)  + \mu^*(A^c \cap E^c)\\
		&> \mu^*(E) + \mu^*(A\cap E^c)  + \mu^*(A^c \cap E^c)\\
		1 - \mu^*(E) &> \mu^*(A\cap E^c)  + \mu^*(A^c \cap E^c)\\
		\mu^*(E^c) &> \mu^*(A\cap E^c)  + \mu^*(A^c \cap E^c)\\
		\end{align*}
		Note that by subadditivity that 
		$		\mu^*(E^c) \leq \mu^*(A\cap E^c)  + \mu^*(A^c \cap E^c)\\
$ Therefore we have a contradiction.  Now assume that $E$ is not measurable.  
	Let the sequence $E_n$ be the sets such that $E \subseteq E_n, E_n = \bigcup_{i=1}^\infty (a_i, b_i], a_i < b_i, \mu^*(E_n) \leq \mu^*(E) + \frac{1}{n}$.  Note we can 
	always choose a cover which overshoots by at most $\frac{1}{n}$, as $\mu^*(E)$ is
	an infimum.  Therefore we can choose $N \in \N$ such that 
	$\frac{1}{n} < \mu^*(E \cap A) + \mu^*(E \cap A^c) - \mu^*(E)$.  Therefore
	$\mu^* \leq \mu^*(E_N) \leq \mu^*(E) + \frac{1}{N} < \mu^*(E \cap A) + \mu^*(E \cap A^c) \leq mu^*(E_N \cap A) + \mu^*(E_N \cap A^c)$, thus we have found a measurable set, $E_N$ such that $\mu(E_n) < \mu^*(E_N \cap A) + \mu^*(E_N \cap A^c)$, and as we have 
	shown previously is a contradiction.  Therefore $A$ must be Lesbesgue measurable.
	\end{itemize}
	\item[5] Let $E \subset [0,1]$ with the property that for some $r \in (0,1)$, 
	every open interval $I$ in $[0,1]$ $\mu^*(E \cap I) \leq r \mu^*(I)$.  
	Note that if we take $I = (0,1)$ we have that $\mu^*(E) = \mu^*(E \cap I) \leq r \mu^*((0,1)) = r$.  Therefore $\mu^*(E) \leq r$.  Since $\mu^*(E)$ is the infimum of 
	the measures of the covers of $E$ with open sets then there exists an open cover
	$\mu^*(\cup_{i=1}^\infty (a_i,b_i)) = r + \epsilon$ for some small $\epsilon > 0$.
	Therefore $\mu^*(E) = \mu^*(E \cap \cup_{i=1}^\infty (a_i,b_i)) 
	\leq \sum_{i=1}^\infty \mu^*(E\cap (a_i,b_i)) \leq r(r+\epsilon)$.  Note that $\epsilon \to 0$, therefore $\mu^*(E) \leq r^2$.  If this holds up to $n$, where 
	$\mu^*(E) \leq r^n$, then there exists an open cover where 
	$E \subseteq \bigcup_{i=1}^\infty (a_i, b_i)$ and $\mu^*(\bigcup_{i=1}^\infty (a_i, b_i)) = r^n + \epsilon$.  Therefore 
	$$
		\mu^*(E) = \mu^*(E \cap 	\bigcup_{i=1}^\infty (a_i, b_i)) 
		\leq \sum_{i=1}^\infty \mu^*(E \cap (a_i, b_i)) \leq 
		\sum_{i=1}^\infty r \mu^*((a_i,b_i)) = r^{n+1} + \epsilon r^n.
	$$
	Since the choice of $\epsilon$ was arbitrary then $\mu^*(E) \leq r^{n+1}$.
	Since $0 < r < 1$, then $\lim r^n = 0$, thus $\mu^*(E) = 0$.  
\end{enumerate}
\end{document}
